\section{Introduction}\label{introduction}

Medical concept normalization --- mapping in-text mentions of
biomedical concepts (diagnoses, medications, labs) to entries in
a standardized ontology or controlled vocabulary such as ICD-10-CM
\autocite{Miftahutdinov2021,Xu2022} --- sits at the foundation of
healthcare data interoperability. Supervised neural classifiers
trained on labeled clinical text have been the state of the art for
most of the last decade, with hierarchical attention over MIMIC
discharge summaries \autocite{Mullenbach2018} and transformer-based
extensions adapting domain-pretrained encoders such as PubMedBERT
\autocite{Gu2021,Huang2022} to the multi-label setting. Two practical
constraints shape how such classifiers can be built and shared in
the current regulatory landscape:

\begin{enumerate}
\def\labelenumi{\arabic{enumi}.}
\tightlist
\item
  \textbf{PhysioNet's responsible-LLM-use policy} \autocite{PhysioNet2025},
  effective 2025-09-24, prohibits sharing credentialed data with
  third parties --- including sending it through commercial LLM APIs.
  This rules out cloud LLM-based classifier construction on
  credentialed corpora and constrains where credentialed data can
  physically reside during training.
\item
  \textbf{Trained-model weights derived from credentialed corpora are
  themselves credentialed-derivative} under PhysioNet's policy, so
  weights cannot be redistributed via standard model-sharing
  channels (HuggingFace Hub, etc.) without going through PhysioNet's
  own MIMIC-IV Models distribution channel.
\end{enumerate}

This technical report documents a fully-local fine-tuning pipeline
for \textbf{PubMedBERT-base} \autocite{Gu2021} on MIMIC-IV-FHIR v2.1
\autocite{Bennett2024,Bennett2023,Johnson2023} restricted to ICD-10-CM
codes in the CMS ACCESS Model FHIR Implementation Guide \autocite{CMS2026}.
The contribution is methodological rather than architectural: every
step from data loading through training through checkpoint
persistence runs on \textbf{consumer Apple Silicon hardware} (M1
MacBook Pro, 16 GB unified memory) via PyTorch's MPS backend,
without ever uploading credentialed content to any cloud service,
training platform, or telemetry endpoint. The same code path
supports CUDA (NVIDIA) and CPU backends; we choose local MPS as
the reference implementation because it is the maximally compliant
configuration and because it demonstrates that research-grade
classifier training on credentialed clinical corpora is reachable
on personal hardware without institutional cloud-compute budgets.

The classifier described here serves as the non-LLM baseline in a
companion paper that benchmarks frontier LLMs against a 6-way
hallucination-aware outcome taxonomy on Synthea-generated
out-of-distribution inputs \autocite{FungPhantomCodes2026}. That paper is
the headline scientific contribution; this report documents the
trained-model arm in full detail so it does not crowd the main
paper's word budget and so future re-implementations have a
self-contained methodology reference.
